\section{Informe del / de la Director/a}

El estudiante de la Carrera de Ingeniería en Informática: \\
\textbf{Ascencio, Felipe Santino}, DNI: 44.670.826, Padrón: 110675 \\
ha cumplido satisfactoriamente con 192 horas de Práctica Profesional Supervisada en 'Empresas de base tecnológica 2' en el marco de su Trabajo Profesional bajo mi supervisión.

El estudiante ha desarrollado satisfactoriamente el Plan de Trabajo oportunamente acordado, por lo que cumplieron con la Práctica Profesional correspondiente, la cual ha servido de base al Trabajo Profesional presentado para su aprobación.

Las tareas asignadas consistieron en participar en el diseño y desarrollo de 'SocioUnido', una plataforma SaaS B2B orientada a modernizar la gestión administrativa, optimizar la recaudación y asegurar el control de accesos en clubes deportivos, asumiendo el rol de responsable de control de calidad, documentación y desarrollo de negocio. El trabajo del estudiante se destacó por su rigurosidad metodológica para planificar y ejecutar pruebas de \textit{software}, su liderazgo en la redacción de la documentación técnica y manuales, y su excelente diseño de estrategias y herramientas orientadas a la comercialización del producto.

El responsable institucional por la organización receptora fue el \textbf{Ing. Franco Di Mauro}, en su rol de \textbf{profesor adjunto} de la materia 'Empresas de base tecnológica 2'. El \textbf{Mg. Ing. Gabriel Piñeiro}, en su carácter de \textbf{director}, estuvo a cargo de la supervisión por parte de la Facultad de Ingeniería.

\vspace{1.5cm}
\noindent\rule{8cm}{0.4pt} \\
FIRMA Y ACLARACIÓN DEL DIRECTOR \\

\vspace{1.5cm}
\noindent\rule{8cm}{0.4pt} \\
FIRMA Y ACLARACIÓN DEL RESPONSABLE DE LA ORGANIZACIÓN RECEPTORA \\

\vspace{0.5cm}

\section{Informe del Estudiante}

Participar en la creación de 'SocioUnido' desde su concepción hasta convertirse en un Producto Mínimo Viable (MVP) funcional ha sido una experiencia sumamente enriquecedora. Construir un producto desde cero, adoptando el verdadero espíritu emprendedor que promueve la materia 'Empresas de base tecnológica 2', me permitió consolidar mis conocimientos técnicos y aplicarlos a una problemática real del mercado. La oportunidad de desarrollar esta Práctica Profesional Supervisada dentro de este marco académico fue vital, ya que aportó una visión integral sobre el ciclo de vida del \textit{software}, la gestión ágil, el trabajo en equipo y la estrategia de negocio, brindándome herramientas invaluables para mi futuro como profesional de la ingeniería.

\vspace{1.5cm}
\noindent\rule{8cm}{0.4pt} \\
FIRMA Y ACLARACIÓN DEL ESTUDIANTE \\
